\documentclass{beamer}

\usetheme{Madrid}
\usecolortheme{spruce}
\usepackage{hyperref}

\title{Bringing TChecker Timed Automata to Life}
\author{TChecker GUI Project}
\date{\today}

\begin{document}

\begin{frame}
  \titlepage
\end{frame}

\begin{frame}{Session Overview}
  \begin{itemize}
    \item Motivation for a visual front-end to TChecker
    \item Timed automata refresher and symbolic reasoning
    \item Demonstrations: Fischer and asynchronous Fischer protocols
    \item Supported verification and simulation workflows
  \end{itemize}
\end{frame}

\begin{frame}{Why a GUI for TChecker?}
  \begin{itemize}
    \item TChecker is a command-line tool for timed automata reachability and liveness analysis that returns textual verdicts only
    \item Timed automata are expressive yet difficult to debug from the command line
    \item Students and engineers need immediate visual feedback on clocks, invariants, and traces
    \item A guided interface lowers the barrier to exploring models while preserving TChecker's rigor
  \end{itemize}
\end{frame}

\begin{frame}{Timed Automata Refresher}
  \begin{itemize}
    \item Following Alur--Dill, a timed automaton is a tuple \(A = (L, \ell_0, C, \Sigma, E, \mathrm{Inv})\)\footnote{\tiny G.~Behrmann, A.~David, and K.~G. Larsen, ``A Tutorial on UPPAAL,'' LNCS 3185, 2004~\cite{behrmann2004uppaal}.}
    \item \(L\): finite set of locations with designated initial location \(\ell_0 \in L\)
    \item \(C\): finite set of real-valued clocks evolving uniformly over time
    \item \(E \subseteq L \times \mathcal{B}(C) \times \Sigma \times 2^C \times L\): edges with guards, events, and reset sets
    \item \(\mathrm{Inv}: L \rightarrow \mathcal{B}(C)\): invariants constraining time at each location
  \end{itemize}
\end{frame}

\begin{frame}{Timed Automata Dynamics}
  \begin{itemize}
    \item \textbf{Guards, actions, invariants}: compared with NFAs, edges carry guards and actions while locations enforce invariants that regulate time
    \item \textbf{Delay transitions}: time may elapse arbitrarily provided the current invariant is maintained throughout the delay
    \item \textbf{Action transitions}: when a guard becomes true, the automaton fires the edge, performs its action, resets clocks, and moves to the target location
  \end{itemize}
\end{frame}

\begin{frame}{Infinite State Motivation}
  \begin{itemize}
    \item \textbf{Infinite concrete space}: clock valuations span \(\mathbb{R}_{\ge 0}\), so concrete states \((\ell, v)\) form an unbounded system
    \item \textbf{Symbolic verification}: we map infinite valuation sets to finite representations that algorithms can manipulate soundly
  \end{itemize}
\end{frame}

\begin{frame}{Symbolic State Exploration}
  \begin{itemize}
    \item TChecker relies on difference bound matrices (DBM) to represent zones\footnote{\tiny TChecker Development Team, ``Advances in Symbolic Exploration for TChecker,'' arXiv:2305.17824, 2023~\cite{herbreteau2023tchecker}.}
    \item Verification reduces to graph search over symbolic states \((\ell, Z)\) with \(Z\) stored as a DBM
    \item Algorithms such as coverage reachability (covreach) explore zones while preserving soundness
  \end{itemize}
\end{frame}

\begin{frame}{Zones and Symbolic Semantics}
  \begin{itemize}
    \item \textbf{Zones}: conjunctions of clock constraints capturing all valuations that satisfy them, enabling finite partitioning of time
    \item \textbf{Symbolic semantics}: states become \((\ell, Z)\); symbolic delay evolves \(Z\), symbolic action checks guards, applies resets, and advances locations
  \end{itemize}
\end{frame}

\begin{frame}{Difference Bound Matrices}
  \begin{itemize}
    \item \textbf{DBM basics}: with \(n\) clocks the matrix has dimension \(n{+}1\), including the zero clock
    \item Entry \(D_{ij}\) stores an upper bound on \(x_j - x_i\), e.g., \((<, 10)\) encodes \(x_j - x_i < 10\)
    \item Closure operations on DBM support efficient zone intersection, extrapolation, and emptiness checks within TChecker
  \end{itemize}
\end{frame}

\begin{frame}{Model Construction Workflow}
  \begin{itemize}
    \item Declare clocks, integer variables, synchronization events, and processes
    \item Configure locations with invariants, committed/urgent flags, and semantic labels
    \item Create transitions with guards, reset sets, and synchronization directions
    \item Export the model to TChecker's \texttt{.tck} format for reproducibility
  \end{itemize}
\end{frame}

\begin{frame}{Interactive Simulation Support}
  \begin{itemize}
    \item Step through time delays and discrete transitions with live DBM snapshots of the current zone
    \item Inspect enabled transitions alongside predicted clock valuations before committing to a move
    \item Replay execution histories to compare schedules and observe how invariants tighten or relax
  \end{itemize}
\end{frame}

\begin{frame}{Verification Features}
  \begin{itemize}
    \item Reachability analysis: find a path to a user-selected label, returning concrete traces
    \item Safety checking: prove that error-labelled states are unreachable or obtain counterexamples
    \item Deadlock detection: confirm every symbolic state admits a successor
    \item Label logic mode: evaluate boolean combinations of labels (e.g., mutual exclusion clauses) under existential or forbidding semantics
  \end{itemize}
\end{frame}

\begin{frame}{Case Study: Fischer Mutual Exclusion}
  \begin{itemize}
    \item Initial step: both processes sit in \texttt{wait} with \(x=0\); after \(x \in (0,1]\) a process writes its id into \texttt{turn} and moves to \texttt{try}
    \item Entry to \texttt{critical} requires guard \(\texttt{turn} = \texttt{pid}\) and invariant \(x \le 2\), preventing simultaneous critical-section access
    \item Simulation view lists the enabled move for each process and displays the DBM so we can watch \(x\) reset on leaving \texttt{critical}
    \item Verification suite runs mutual exclusion on \texttt{critical\_P1}, \texttt{critical\_P2} and a safety query ensuring the synthetic \texttt{Bad} label is unreachable
  \end{itemize}
\end{frame}

\begin{frame}{Case Study: Asynchronous Fischer Variant}
  \begin{itemize}
    \item Offsets introduce distinct clocks \(x_1, x_2\); each process must wait its own guard window before claiming \texttt{turn}
    \item The first enabled step shows only one process eligible to proceed while the other remains delayed, reflected in the DBM boundaries
    \item Subsequent transitions highlight how asynchronous retries reshape the interleaving yet still enforce the \texttt{turn} equality before \texttt{critical}
    \item Verification compares deadlock freedom and mutual exclusion between the synchronous and asynchronous models to quantify timing slack
  \end{itemize}
\end{frame}

\begin{frame}{Demo Script Summary}
  \begin{enumerate}
    \item Import the Fischer model (\texttt{fischer.tck}) and inspect declarations
    \item Step through the simulator to observe enabled transitions and clock constraints
    \item Run mutual exclusion and safety verification; review returned traces
    \item Switch to the asynchronous variant (\texttt{fischer-async.tck}) and repeat analysis
    \item Compare results to highlight the impact of altered timing assumptions
  \end{enumerate}
\end{frame}

\begin{frame}{Key Takeaways}
  \begin{itemize}
    \item Timed automata modeling benefits from immediate visualization of zones and traces
    \item Integrated simulation plus verification shortens the feedback loop for debugging
    \item Fischer examples showcase how timing parameters influence mutual exclusion guarantees
    \item The GUI encapsulates TChecker's verification strengths within an accessible workflow
  \end{itemize}
\end{frame}

\begin{frame}{Next Steps}
  \begin{itemize}
    \item Iterate on zone-based BFS heuristics to accelerate large-model exploration
    \item Add full support for diagonal constraints through refined DBM normalisation and pruning
    \item Profile verification runs on benchmark suites to validate performance gains
  \end{itemize}
\end{frame}

\begin{frame}{Questions}
  \centering
  Thank you for your attention!\\
  We welcome discussion on new verification scenarios and benchmarks.
\end{frame}

\begin{frame}{References}
  \begin{thebibliography}{9}
    \bibitem{behrmann2004uppaal}
      G.~Behrmann, A.~David, and K.~G. Larsen.
      \newblock A Tutorial on UPPAAL.
      \newblock In \emph{Formal Methods for the Design of Real-Time Systems}, LNCS~3185, 2004.
      \newblock \url{https://www.seas.upenn.edu/~lee/09cis480/papers/by-lncs04.pdf}
    \bibitem{herbreteau2023tchecker}
      TChecker Development Team.
      \newblock Advances in Symbolic Exploration for TChecker.
      \newblock Technical report, arXiv:2305.17824, 2023.
      \newblock \url{https://arxiv.org/pdf/2305.17824}
  \end{thebibliography}
\end{frame}

\end{document}
